\section{Support induction and Chow vanishing}
\label{sec:support}

We now combine the arbitrary-color bound of
Lemma~\ref{lem:large-vertex-test} with the compact lifts of
Proposition~\ref{prop:controlled-lifts}.  We will make successive parts
of the complete product of three-color entries vanish over dense opens
of its base support, until no support remains.  Throughout this section,
\[
 g=10^{60},\qquad D_*=\binom g3,\qquad n=3D_*,
 \qquad U_0=\Q^g.
\]
The basis of \(U_0\) is the set of colors.  An \emph{entry} is one factor
with three specified color inputs.  Initially there is one entry for each
increasing triple of distinct basis colors.  Products and their
factorizations are taken in the sign-coinvariant operation algebra of the
weighted spaces from the preceding section.

The graph charts, their compact heavy bases, and their smooth-core opens
will be denoted by
\[
 X^\Gamma(k),\qquad \mathcal B_\Gamma,\qquad
 \mathcal S_\Gamma\subset\mathcal B_\Gamma,
\]
respectively.  Every graph in this section has unmarked type \((g,0)\);
its vertex flags are heavy marks.  An operation supported on a vertex is
zero on assignment components on which its active labels are not all at
that vertex.

\subsection{Advancing support with an arbitrary multiplier}

Boundary-support methods also appear in tautological vanishing results
such as \cite[Theorem~1.1]{GraberVakil}, where support is measured by
rational components. Here we lower the dimension of the image in the
heavy base, using operational vanishing over smooth test families.
We first make this induction precise.  A term in it consists of
active operations, pulled to disjoint sets of light labels, acting on a
cycle on the full labeled chart.  The cycle need not factor over the
labels or over the vertices.

\begin{lemma}[Support advancement]
\label{lem:support-advancement}
Let a cycle multiplier on \(X^\Gamma(n)\) have base image contained in an
irreducible closed subspace \(Z\subset\mathcal B_\Gamma\).  After passing
to a refinement if necessary, suppose that its generic base image lies
in \(\mathcal S_\Gamma\).  Let \(L\) be a factor involving a specified
subset of the active entries.

Suppose that the pull of \(L\) vanishes operationally over a dense open
of a smooth test family dominating \(Z\cap\mathcal S_\Gamma\).  Then its
action can be written as a finite sum of cycles with strictly smaller
base-image dimension.  Only the entries in \(L\) are consumed.

More precisely, take a smooth test base \(S\), proper and generically
finite over \(Z\cap\mathcal S_\Gamma\).  If the vanishing holds off base
codimension at least \(a\), the terms remaining over \(\mathcal S_\Gamma\)
have base-image dimension at most \(\dim Z-a\).  The other terms lie over
\(Z\setminus\mathcal S_\Gamma\) and require a strict graph refinement.
If one first replaces \(S\) by the total space of one of its curve
families to obtain a section, the corresponding bound is
\(\dim Z-a+1\).
\end{lemma}

\begin{proof}
The proper cut-chart maps cover the preimages of boundary strata.
Consequently a cycle whose generic heavy curve has a finer dual graph
can be lifted to a chart of that graph.  These lifts do not increase the
dimension of its base image: heavy gluing is finite.

We use the following elementary property of rational Chow groups.  If
\(T'\to T\) is proper and surjective, every integral cycle on \(T\) has,
after multiplication by a nonzero rational number, a proper lift
generically finite over it.  Indeed, choose a closed point of the fiber
over its generic point, take its finite residue-field extension, and
close up that point.  The pushforward is the original cycle multiplied
by the generic degree.  The same property holds for the finite-type Deligne--Mumford charts
here with rational coefficients \cite[Proposition~B.19]{BaeSchmitt}.

For the test base, take a finite level cover of the smooth heavy-curve
stacks and resolve the appropriate component over
\(Z\cap\mathcal S_\Gamma\)
\cite[Theorems~1.0.2 and~1.0.3]{Wlodarczyk}.  This gives a smooth quasi-projective scheme
proper and generically finite over that space.  A full level of order at least three removes automorphisms in positive
genus \cite[Sections~5.3--5.4 and~5.14]{DeligneMumford},
\cite[IV.2.7--2.9 and IV.3.1--3.2]{DeligneRapoport}; the genus-one
vertices here have flags, and smooth genus-zero factors are already
schemes.  Lift the cycle multiplier properly to the resulting family by
the preceding argument.

Let \(A\subset S\) have codimension at least \(a\), and suppose that \(L\)
is zero operationally over \(S\setminus A\).  Its action on the lifted
multiplier restricts to zero there.  Chow localization expresses the
result as a pushforward of cycles lying over \(A\).  Projection formula
and proper pushforward return this equality to the original family.
Because \(\dim S=\dim Z\), these cycles have base-image dimension at most
\(\dim Z-a\).  Closing up inside the inverse image of \(Z\) can introduce
additional cycles only over \(Z\setminus\mathcal S_\Gamma\), a proper
closed subset of \(Z\).  Refining their generic dual graphs places them
on the required cut charts.

To obtain a section of a smooth curve family \(C\to S\), pull that
family back along \(C\to S\), where it has its diagonal section.
The new test base is smooth and has dimension \(\dim Z+1\).
Although this base change is not generically finite, the cycle multiplier
still admits a proper lift generically finite over itself: apply the
same generic-point argument to its pullback.  A support of codimension
at least \(a\) in the new base has image dimension at most
\(\dim Z+1-a\), proving the last assertion.

These operations concern the chosen factor \(L\).  All remaining
operations are carried along by pullback and projection formula and can
be applied afterwards to the resulting cycle multipliers.
\end{proof}

The operational hypothesis in this lemma is essential.  In its
applications below, it follows from identities over smooth test
families, rather than from identities separately on geometric fibers.
For a pure operation, Proposition~\ref{prop:controlled-lifts} identifies
its restriction with a pullback from a smooth curve power.  Vanishing
of its Chow class on that power is therefore operational vanishing and
can be pulled further to any cycle multiplier.  The formation of the
projector and the projected diagonal commutes with these base changes:
the relative diagonals and the integrations along smooth proper curve
powers have the requisite flat or refined-Gysin base-change identities.
For residual operations, the explicit base-divisor factors in
Proposition~\ref{prop:controlled-lifts} give the required open vanishing.
The same is true of the controlled singleton relations below.

There is also no restriction on the multiplier hidden in our use of
sign coinvariants.  In characteristic zero, signed averaging identifies
coinvariants with the corresponding signed invariant subspace.  A
factorization \(L*H\) of operations becomes a sum over shuffles of
products of the separately averaged \(L\) and \(H\), pulled to disjoint
slot sets.  In each summand, apply \(L\) first to the multiplier.
The factor \(H\) remains an averaged operation on its own active slots;
it is pulled equivariantly through later cuts before it acts.
This proves that the lemma applies repeatedly, even after the multiplier
has become an arbitrary cycle produced by an earlier step.
More explicitly, every later factorization of \(H\) is first an equality
of its signed-average representatives on its own active slots.  Pull
that equality to the full chart and only then cap with the multiplier.
No permutation is applied to the multiplier, and no invariance of the
multiplier is required.  Restriction to an assignment component or
pullback through a further cut preserves this equality.

Initially the heavy base has dimension \(3g-3\).  Every advancement
strictly lowers the base-image dimension.  Thus on any branch of the
finite expansions there can be at most
\begin{equation}
\label{eq:advancement-number}
 3g-2
\end{equation}
advancements before the support is empty.

\subsection{The elementary vanishing tests}

On a current chart, Proposition~\ref{prop:controlled-lifts} expands an
ancestor entry into pure entries at current vertices, products of a pair
operation with a controlled singleton, and residual cubic operations.
Call the pair-times-singleton terms \emph{noise entries}.
Perform these expansions one entry at a time.  If a residual summand is
chosen, apply Lemma~\ref{lem:support-advancement} to it immediately, at
the cost of that one entry.  We therefore never retain residual entries
while counting the pure and noise entries.

The pure entries at any one vertex always use the same chosen symmetric
class \(G_{h,B}\).  This is important both for the finite coefficient
bound and for the common-color test below.

\begin{lemma}[Noise and small-genus tests]
\label{lem:elementary-entry-tests}
Set
\[
 R_\Gamma=2|V(\Gamma)|+2|E(\Gamma)|.
\]
If \(R_\Gamma+1\) noise entries have their singleton at the same color,
their product advances support while consuming those entries.
If a genus-\(h\) vertex has more than \(\binom{2h}{3}\) pure entries,
the product of any \(\binom{2h}{3}+1\) of them advances support.
\end{lemma}

\begin{proof}
Over smooth generic core data, the one-light space at a vertex is its
curve.  The permitted singleton classes lie in the span of its unit,
canonical class, and flag points.  Over all vertices this gives a space
of dimension at most
\[
 \sum_v(2+|B_v|)=2|V|+2|E|=R_\Gamma.
\]
A one-slot operation is odd in the sign-coinvariant algebra.
At a fixed color, a product of \(R_\Gamma+1\) such elements is zero by
multilinearity and alternation.  The accompanying pair operations have
even arity and can be grouped separately.  This proves the first
vanishing, and it holds on an open of each smooth test family by the
controlled singleton descriptions and Chow localization.

For the second assertion apply Proposition~\ref{prop:tensor-test}.
Every tested pure entry is a value of the same odd-coefficient form
\(\varphi\in\mathfrak B\otimes\bigwedge^3 W_h^\vee\).
Its values belong to the span of at most \(\binom{2h}{3}\) odd
coefficients.  A longer product vanishes by supercommutativity.
The faithful same-map tensor test proves vanishing of the operation
modulo positive base support.  Only finitely many coefficient tests are
needed in a fixed arity, so their supports can be combined without
losing positive codimension.  Apply
Lemma~\ref{lem:support-advancement}.
\end{proof}

\subsection{A repeated-color test at a flagged vertex}

The surface test of Lemma~\ref{lem:large-vertex-test} is needed only
while one vertex has almost all the genus.  On every other chart a less expensive advancement follows
from the flag-centered curve algebra.

\begin{lemma}[Repeated-color advancement]
\label{lem:repeated-color}
Suppose that a vertex of genus \(h\) has a flag and that
\(k>\binom{h+1}{2}\) pure entries contain one common color \(i\).
Their product is a finite linear combination of products \(L*H\) with
the following properties:
\begin{enumerate}[label=\textup{(\roman*)}]
\item \(L\) involves at most \(h+2\) entries and vanishes modulo positive
base support;
\item \(H\) is still a product of individual pure entries at that vertex;
\item before removing the entries in \(L\), each term has the same color
multiset as the original product.
\end{enumerate}
Consequently each term admits a support advancement consuming at most
\(h+2\) entries.
\end{lemma}

\begin{proof}
A symmetric cubic class defines an alternating three-color operation.
Put the common color first in each entry.  The product, as a function of
the other \(2k\) inputs, factors through
\[
 \bigwedge^k\bigl(\bigwedge^2 U_0\bigr).
\]
This holds already on the compact chart: alternation inside a pair comes
from cubic symmetry, and interchange of two cubic entries has sign
minus because their arity is odd.

Let \(p\in\Q[\mathfrak S_{2k}]\) be the corresponding projector on the
input slots.  The exterior functor displayed above is zero on a vector
space of dimension \(h+1\).  By the Schur--Weyl decomposition and its row cutoff
\cite[Theorems~5.12.2, 5.18.4, and~5.22.1]{EtingofEtAl2011},
\(p\) lies in
the two-sided ideal generated by the alternator on \(h+2\) slots:
the quotient of the group algebra by that ideal acts faithfully on
tensor powers in dimension \(h+1\).  Indeed, its simple blocks are
indexed by partitions; precisely those with at most \(h+1\) rows
survive in that tensor power, and the \((h+2)\)-alternator generates
all the other blocks.  Thus \(p\) is a finite sum of
operators of the form
\[
 c\,\sigma\,\operatorname{Alt}_{h+2}\,\tau,
 \qquad c\in\Q,\quad \sigma,\tau\in\mathfrak S_{2k}.
\]
In one such term, first permute the color inputs by \(\tau\), and
interpret \(\sigma\) as their assignment to the pairs of entries.
Only entries meeting the \(h+2\) alternated positions vary during the
alternation.  There are at most \(h+2\) of them.  Group their alternated
product into \(L\), leaving the other entries in \(H\).
This proves the asserted factorization and the color-multiset statement.

It remains to prove the vanishing of \(L\).  Work on a smooth test base
modulo positive base support and center the curve projector at the chosen
flag section.  By Lemma~\ref{lem:point-centered-algebra}, its class \(e\)
has
\[
 \int e=1,\qquad \int e^2=0.
\]
Evaluation at the section is pairing with \(e\), is multiplicative, and
kills the pure odd space \(W_h\).  Hence for standard odd vectors \(u,v\),
their product in the coefficient Frobenius algebra has the form
\begin{equation}
\label{eq:point-centered-pair}
 uv=\omega(u,v)e+(uv)_{\mathrm{pure}}.
\end{equation}
There is no unit component, since its coefficient is the value of \(uv\)
at the section.

Test the colors by \(z:U_0\to W_h\), and write \(a=z(i)\).  Each entry
becomes
\[
 T_a(u,v)=\varphi(a,u,v),
\]
an alternating two-form with odd coefficients.  Contracting a variable
slot between two such entries gives zero.  To see this, the full
Frobenius contraction is, up to its fixed Koszul sign,
\[
 \int (au)(av)=0,
\]
because \(a\) is odd and \(a^2=0\).
The removed hyperbolic contraction in
\eqref{eq:point-centered-pair} is a multiple of \(\int e^2=0\);
orthogonality eliminates its mixed contractions.
The remaining symplectic contraction of the pure products is exactly
the stated contraction of the two cubic tensors.

An internal contraction of the two variable slots in a single \(T_a\)
is one and the same odd scalar, say \(\ell(a)\), up to the orientation
sign.  Two internal contractions in distinct entries therefore give
zero, because \(\ell(a)^2=0\).

After alternation, \(L\) is an ordinary alternating
\((h+2)\)-form on \(W_h\), with coefficients in the supercommutative
algebra.  Every summand of its double symplectic contraction is zero:
if one contraction joins distinct entries, perform that external
contraction first; otherwise the two contractions are internal in
distinct entries and give the square just considered.  All other
variables can be held arbitrary in the external identity, so this
reasoning also covers chains of contractions.

For \(h\ge2\), double contraction
\[
 \bigwedge^{h+2}W_h^\vee
 \longrightarrow \bigwedge^{h-2}W_h^\vee
\]
is injective by symplectic Lefschetz decomposition.  For \(h=0,1\),
the source alternation is already zero.  Thus every color test of \(L\)
vanishes.  Proposition~\ref{prop:tensor-test} and
Lemma~\ref{lem:support-advancement} finish the proof.
\end{proof}

\subsection{The finite color budget}

Set
\begin{equation}
\label{eq:color-constants}
 K=10^5,\qquad m_0=2000g^{2/3},\qquad
 C_K=\binom{2K}{3}.
\end{equation}
These are integers.  Reserve \(4K+20\) disjoint sets of \(m_0\) basis
colors each, and let \(S_0\) be all the remaining colors.  This is possible
because
\[
 \frac{(4K+20)m_0}{g}
   =8.0004\cdot10^{-12}<0.1;
 \qquad |S_0|>0.9g.
\]
For a current graph put
\[
 b=\max_v h_v,\qquad l=g-b.
\]
Under refinement the maximum vertex genus cannot increase, so \(l\)
cannot decrease.  The number of edges cannot decrease either.
Unmarked stability gives
\begin{equation}
\label{eq:global-noise-bound}
 |V|\le2g-2,\qquad |E|\le3g-3,\qquad R_\Gamma\le10g.
\end{equation}

An advancement by a residual entry, by a noise test, by a small-genus
test, or by Lemma~\ref{lem:repeated-color} will be called
\emph{ordinary}.  These are descriptions of the cost in entries, not
additional geometric assumptions.

The two regimes below use different information about the remaining
colors.  While \(l<K\), we remove whole entries without changing the
three colors of any retained entry; the reserved blocks can therefore
be counted by their original triples.  We use the repeated-color
rewrites only once \(l\ge K\).  They can change which colors occur
together in an entry, but preserve the total multiset of color
occurrences before the advancing factor is removed.  In that regime
we count occurrences in \(S_0\).  Since \(l\) cannot decrease, we never
need an original-triple count after making such a rewrite.

\paragraph{Small genus deficit: \(l<K\).}
There is a unique vertex of genus \(b\), since \(b>g-K>g/2\).
The total genus of all the others is at most \(l\), and stability gives
\begin{equation}
\label{eq:small-graph-size}
 |V|-1\le2l,\qquad |E|\le3l,\qquad R_\Gamma\le10K+2.
\end{equation}
For completeness, sum \(2h_v-2+\operatorname{val}(v)\ge1\) over all
vertices except the large one.  The sum equals
\(2l-\operatorname{val}(v_{\mathrm{large}})\), proving the first
inequality.  The genus formula then gives the second.

After expanding the active entries, use a residual immediately.
If a color occurs as a singleton in more than \(R_\Gamma\) noise entries,
use Lemma~\ref{lem:elementary-entry-tests}.  Otherwise there are at most
\(g(10K+2)\) noise entries.  If a small vertex has more than \(C_K\)
pure entries, use the same lemma there.  Otherwise the number of pure
entries off the large vertex is at most \(2K C_K\).
Each ordinary advancement in this regime costs at most
\[
 M_K=\max\{10K+3,C_K+1\}=C_K+1.
\]
As specified above, retained entries in this regime still have their
original triples, including those internal to each unused reserved
block.

Suppose none of these ordinary tests applies.  Choose a reserved color
block \(E_0\) that has not been used for a large-vertex test.  Let
\(D_0=\binom{m_0}{3}\).  Among its original \(D_0\) entirely internal
entries, the number that are no longer active pure entries at the large
vertex is at most
\begin{equation}
\label{eq:small-loss}
 3gM_K+g(10K+2)+2K C_K
 <10^{17}g<0.01D_0.
\end{equation}
Here \eqref{eq:advancement-number} bounds all preceding ordinary losses.
Earlier large-vertex tests used only entries internal to different,
disjoint reserved blocks, so they removed none of these entries.
The numerical inequalities follow from
\[
 C_K<1.334\cdot10^{15},
 \qquad D_0>10^9g^2,
 \qquad g=10^{60}.
\]
Thus more than \(0.99D_0\) internal entries remain pure at the large
vertex.

Consume the product of all those entries.  On a smooth test base,
obtain a section if necessary as in
Lemma~\ref{lem:support-advancement}.  Its dimension is at most \(3g-2\);
since \(b>g-K\),
\[
 3g-2<(b-1)^2-1.
\]
We may therefore use Lemma~\ref{lem:large-vertex-test}.
Choose an even integer \(e\) between \(m_0/2-2\) and \(m_0/2\).
For \(d_e=\binom e3\) and \(t_b=25(b-1)^2-1\), the explicit bounds are
\[
 d_e>D_0/10,\qquad t_b<25g^2,\qquad t_b<d_e/2,
\]
and consequently
\begin{equation}
\label{eq:large-test-threshold}
 D_0-d_e/2+t_b
   <(0.95+2.5\cdot10^{-8})D_0<0.99D_0.
\end{equation}
For the first bound, \(m_0\ge128\) and
\[
 \frac{\binom e3}{\binom{m_0}3}
 \ge \left(\frac12-\frac4{m_0}\right)^3>\frac1{10}.
\]
The selected product thus vanishes modulo base support of codimension
at least \(b-1\).  Applying
Lemma~\ref{lem:support-advancement}, it either lowers the base-image
dimension by at least \(b-2>g-K-2\), or gives a boundary remainder
requiring a strict refinement.  Call this an \emph{expensive}
advancement and mark \(E_0\) as used.

There are enough fresh blocks for every such advancement.
Those that require a refinement start with fewer than \(3K\) edges by
\eqref{eq:small-graph-size}; along a branch there are at most \(3K+1\)
of them.  Every other expensive advancement lowers the base-image
dimension by at least \(g-K-2\).  There can be at most three such
advancements with nonempty resulting support, because
\[
 4(g-K-2)>3g-3.
\]
A fourth can be terminal, making the support empty.
In total at most \(3K+5<4K+20\) blocks are needed.
No expensive advancement consumes an occurrence of a color in \(S_0\).

\paragraph{Large genus deficit: \(l\ge K\).}
Once this regime is reached it persists.  We continue to expand
ancestors using the same pure lift at each current vertex and peel
residuals immediately.  A noise test consumes at most \(10g+1\) entries.
If no such test applies, at most \(10g^2\) noise entries remain.

All ordinary advancements on the entire preceding path, including
those in the small-deficit regime, cost at most \(11g\) entries each:
\[
 C_K+1<11g,\qquad 10g+1<11g,\qquad h_v+2<11g.
\]
There have been at most \(3g-2\) advancements in total.
Each consumed entry has three color occurrences.  The cumulative
ordinary loss of occurrences in \(S_0\) is therefore less than
\(99g^2\).  The unpeeled noise accounts for at most another \(30g^2\).
Expensive advancements used no \(S_0\) colors.  Hence the number of
pure occurrences of colors from \(S_0\) is at least
\begin{equation}
\label{eq:remaining-pure-occurrences}
 |S_0|\binom{g-1}{2}-140g^2.
\end{equation}
The allowance \(140g^2\) is larger than the sum \(129g^2\) of the
two losses just obtained.
The permutations used in Lemma~\ref{lem:repeated-color} preserve the
color multiset before its advancing factor is removed, so the same
bound remains valid after any number of those rewrites.

Discard any identically zero product, in particular a product containing
a pure entry with a repeated color.  Suppose that every color
\(i\in S_0\), at every vertex \(v\), belongs to at most
\(\binom{h_v+1}{2}\) pure entries.  The total number of pure occurrences
would then be at most
\[
 |S_0|\sum_v\binom{h_v+1}{2}
 \le \frac{|S_0|}{2}(b+1)\sum_v h_v
 \le \frac{|S_0|g(b+1)}2.
\]
This contradicts \eqref{eq:remaining-pure-occurrences}.  Indeed, the
lower bound minus this upper bound is
\[
 \frac{|S_0|}{2}\bigl(g(l-4)+2\bigr)-140g^2
 >\bigl(0.45(K-4)-140\bigr)g^2>0.
\]
There is therefore a color and vertex to which
Lemma~\ref{lem:repeated-color} applies.
Every vertex in this regime has a flag: a connected graph without any
flags at one of its vertices could only be the one-vertex, no-edge
graph, for which \(l=0\).
The lemma supplies another ordinary advancement, at the claimed cost.

\subsection{Completion}

\begin{theorem}
\label{thm:chow-vanishing}
The class \(Y\in\Salg_{g,n}\) constructed above satisfies
\[
 q_{g,n}(Y)=0
 \quad\text{in }\CH^*(\Mbar_{g,n};\Q).
\]
\end{theorem}

\begin{proof}
Begin on \(X_{g,\varnothing}(n)\) with the product of the cubic
operations \(F^0_{g,\varnothing}\), one for each increasing color triple,
and with the fundamental cycle as multiplier.  Apply
Proposition~\ref{prop:controlled-lifts} and the preceding induction.
At every stage, each nonzero branch admits an advancement: use a
residual or noise test when available, the small-vertex and reserved-block
tests if \(l<K\), and otherwise the repeated-color test.
The counts show that all required active entries and unused blocks are
available at the moment they are needed.

Each individual expansion and localization involves finitely many
cycles.  Each branch has length at most \(3g-2\), by
\eqref{eq:advancement-number}.  Thus the process ends with empty
support and proves that the full product is zero in the sign-coinvariant
Chow group of the weighted space.

The weighted-to-unweighted comparison in
Lemma~\ref{lem:weighted-pullback} identifies the pullback of the
raw cubic with \(q(F_J^{\mathrm{st}})\).  After signed averaging, the
coefficient of the ordered color word \(\mathbf i\) is a nonzero
rational scalar times
\[
 q_{g,n}\!\left(
 \sum_{\tau:\,\tau\mathbf i=\mathbf i}
   \operatorname{sgn}(\tau)\,
   \tau\!\left(\prod_J F_J^{\mathrm{st}}\right)
 \right)=q_{g,n}(Y).
\]
Pulling back the zero weighted class proves the theorem.
\end{proof}

\begin{proof}[Proof of Theorem~\ref{thm:main}]
Theorem~\ref{thm:chow-vanishing} places \(Y\) in the Chow kernel.
Proposition~\ref{prop:detector} supplies a linear map annihilating the
specified Pixton span and taking a nonzero value on \(Y\).  Thus
\[
 Y\in\ker q_{g,n}\setminus\Pspan_{g,n}.
\]
The stable pair \(g=10^{60}\), \(n=3\binom g3\) is a counterexample to
the proposed equality for all stable pairs.
\end{proof}
